\section{Preliminaries}\label{sec:prelim}

We collect the definitions and conventions used throughout the paper.
Standard references for universal algebra and equational
logic include Burris and Sankappanavar~\cite{burris1981course},
McKenzie, McNulty, and Taylor~\cite{mckenzie1987algebras},
and Taylor~\cite{taylor1979equational}.

\subsection{Magmas and equational laws}

\begin{definition}\label{def:magma}
A \emph{magma} is a pair $(M, *)$ where $M$ is a nonempty set and
$* \colon M \times M \to M$ is a binary operation.  No axioms are
imposed on~$*$.
\end{definition}

We work in the first-order language $\mathcal{L} = \{*\}$ with a
single binary function symbol.  \emph{Terms} are built from a countable
set of variables $\{x, y, z, w, \ldots\}$ and the operation~$*$
according to the usual recursive definition: each variable is a term,
and if $s$ and~$t$ are terms then $s * t$ is a term.

\begin{definition}\label{def:eq-law}
An \emph{equational law} is a formal identity $s = t$ where $s$
and~$t$ are terms in~$\mathcal{L}$.  A magma $(M,*)$ \emph{satisfies}
$s = t$ if for every assignment $\sigma$ of the variables to elements
of~$M$, the identity $\sigma(s) = \sigma(t)$ holds in $(M,*)$.
\end{definition}

\begin{definition}\label{def:implication}
Let $E_1$ and $E_2$ be equational laws.  We say that $E_1$
\emph{implies} $E_2$, written $E_1 \models E_2$, if every magma
satisfying $E_1$ also satisfies~$E_2$.
\end{definition}

By Birkhoff's theorem~\cite{birkhoff1935structure}, $E_1 \models E_2$
if and only if $E_2$ can be derived from $E_1$ using the rules of
equational logic (reflexivity, symmetry, transitivity, substitution,
and congruence).

\subsection{The equation corpus}

Following Bolan et al.~\cite{bolan2025equational}, we consider the
finite set $\Eq$ of all equational laws with at most four occurrences
of~$*$.  This set has $|\Eq| = 4694$ elements, partitioned into $1415$
equivalence classes under mutual implication.  Of the
$4694^2 = 22{,}033{,}636$ ordered pairs, $8{,}178{,}279$ (approximately
$37.12\%$) are true implications.

\subsection{Structural features of equational laws}\label{subsec:features}

We define several syntactic features of an equational law $E\colon s = t$
that are used in our prediction rules.

\begin{definition}\label{def:structural}
Let $E\colon s = t$ be an equational law.
\begin{enumerate}[label=(\roman*)]
\item The \emph{variable count} $\var(E)$ is the number of distinct
  variables appearing in~$E$.
\item The \emph{size} $\siz(E)$ is the total number of occurrences
  of~$*$ in both $s$ and~$t$.
\item The \emph{signature} of~$E$ is the pair $(a,b)$ where $a$
  (resp.~$b$) is the number of occurrences of~$*$ on the left-hand
  side (resp.\ right-hand side).
\item We say $E$ is \emph{bare} if one side of $E$ is a single
  variable.  Equivalently, $E$ has signature $(0, b)$ or $(a, 0)$
  for some $a$ or~$b$.
\item When $E$ is bare with a single variable $x$ on one side, we say
  $x$ is the \emph{anchor variable}.  We say $E$ is \emph{lhs-var}
  if $x$ appears on the left: $x = t$.
\end{enumerate}
\end{definition}

Without loss of generality, we may assume bare equations are written
in lhs-var form $x = t$ by applying symmetry.

\begin{definition}\label{def:projections}
Let $E\colon x = t$ be a bare, lhs-var equational law.
\begin{enumerate}[label=(\roman*)]
\item We say $E$ has the \emph{left-occurrence} property,
  $\Lx(E) = \textsc{true}$, if the anchor variable~$x$ appears
  as the left child of the root of~$t$.
\item We say $E$ has the \emph{right-occurrence} property,
  $\Rx(E) = \textsc{true}$, if the anchor variable~$x$ appears
  as the right child of the root of~$t$.
\item We say $E$ has the \emph{right-projection} property,
  $\RP(E) = \textsc{true}$, if the right-hand side~$t$ can be
  written as $u * x$ for some term~$u$ not containing~$x$.
\end{enumerate}
\end{definition}

\begin{definition}\label{def:topshape}
The \emph{top shape} of a term $t$ records the structure of its
two outermost levels.  If $t = s_1 * s_2$ and each $s_i$ is either
a variable (denoted $v$) or a product $s_i = u * w$ (denoted $m$),
then the top shape of~$t$ is one of $v$-$v$, $v$-$m$, $m$-$v$,
or $m$-$m$.  If $t$ is a single variable, its top shape is~$v$.
\end{definition}

\subsection{Projection algebras}

\begin{definition}\label{def:proj-alg}
A magma $(M, *)$ is a \emph{left-projection algebra} if $a * b = a$
for all $a, b \in M$.  It is a \emph{right-projection algebra} if
$a * b = b$ for all $a, b \in M$.
\end{definition}

Left- and right-projection algebras are among the simplest magmas.
Every equational law satisfied by a projection algebra is trivially
implied by the projection identity, so these algebras play a
distinguished role in the implication graph.
